\documentclass[11pt,journal,onecolumn]{IEEEtran} 
\pagenumbering{gobble}
\begin{document}


\title{Title}


\author{\IEEEauthorblockN{Author 1\IEEEauthorrefmark{1} and Author 2\IEEEauthorrefmark{2}}

%% If more than two authors
%\author{\IEEEauthorblockN{Author 1\IEEEauthorrefmark{1}, Author 2\IEEEauthorrefmark{2}, and Author 3\IEEEauthorrefmark{3}}

\IEEEauthorblockA{\IEEEauthorrefmark{1}Johns Hopkins University Applied Physics Lab, USA, email: x@y.edu}

\IEEEauthorblockA{\IEEEauthorrefmark{2}Princeton University, USA, email: x@y.edu}

%% Affiliation of third author
%\IEEEauthorblockA{\IEEEauthorrefmark{3}University, USA, email: x@y.edu}

}

\maketitle


\begin{abstract}
The emergence of quantum computing provides us the possibility of solving tasks that might take years classically in just a few minutes. For certain problems, quantum computing exhibits quantum supremacy, meaning that the quantum solution runs exponentially faster than classical algorithms and is able to completely take over classical computers. This high efficiency of quantum computing comes not only from the hardware but also the software, quantum algorithms. The algorithms utilize the qubits to make calculations in order to fulfill specific tasks with the lowest time complexity possible. One such algorithm is named the Grover's algorithm, which is able to perform database search in $\mathcal{O}(\sqrt{N})$, and it runs much faster than the traditional algorithm that takes $\mathcal{O}(N)$ time to solve the same task. For example, when the task is to find the even integers from $N$ integers, traditional computation will need to run through all of the $N$ integers one by one, making at least $N$ steps of calculation, while by using Grover's algorithm only around $\sqrt{N}$ calculations are needed. This exponential speed-up makes Grover's algorithm one of the most important quantum algorithms. Grover's algorithm has a wide application in many fields and is able to improve the time complexity exponentially. One task that can be solved using Grover's algorithm is the satisfiability problem. This type of problem asks the computer to find a set of values (commonly true or false) for several variables such that they satisfy certain constraints. We use k-SAT problems to refer to satisfiability problems with $k$ boolean variables to be determined. Grover's algorithm can effectively solve the k-SAT problem by performing the database search on $2^N$ possible states of the variables. The algorithm's square root optimization on searching helps to improve the efficiency of this solution significantly. Furthermore, this optimization of Grover's algorithm may play a more important role when k grows larger, and consequently the efficiency of the quantum solution could improve faster relative to the traditional solution. Yet this hypothesis is never tested due to the lack of a general k-SAT quantum algorithm. No quantum algorithms solving k-SAT problems where k is greater than 3 have been proposed, thus no test has been performed to compare the quantum solution and the classical solution on more general k-SAT problems. In this research, we formulate a general quantum solution for k-SAT problem and compare such solution with the best classical algorithm to determine whether and when the quantum algorithm performs better on satisfiability problems. The comparison will be done through both theoretical deduction as well as real-world implementation. At the end of this research, we will determine whether the proposed quantum algorithm outperforms the classical algorithm on solving $k$-satisfiability problems.
\end{abstract}

\end{document}